% =====================================================================
%  5.9 RESUMEN DE LA SOLUCION PROPUESTA
%  Incluye OBLIGATORIAMENTE la metodologia empleada.
%  Rubrica: para "Excelente" hay que justificar la solucion de forma
%  detallada y extensa, e incluir la metodologia JUSTIFICANDO SU EMPLEO.
% =====================================================================

\chapter{Solución propuesta y metodología}
\label{cap:solucion}

\section{Visión general de la solución}

La solución es una aplicación web de apoyo a la decisión clínica cuyo núcleo
es un motor de evaluación: las reglas STOPP/START v3 se mantienen aparte del
mecanismo que las aplica, de modo que los criterios pueden actualizarse sin
reescribir la lógica de la aplicación. El detalle interno del motor se
desarrolla en los capítulos de arquitectura y diseño.

El profesional recorre un flujo en dos pestañas. En medicamentos selecciona
los fármacos del paciente; el motor evalúa el caso y marca cada criterio como
activo o inactivo. En diagnóstico se completa la comorbilidad y el motor
vuelve a evaluar, con la misma distinción entre criterio activo e inactivo.
El resultado final agrupa los avisos STOPP (prescripción potencialmente
inapropiada) y START (posible omisión de tratamiento indicado).
La~\cref{fig:flujo-vision} resume este recorrido. La exportación e importación
del caso se describen en capítulos posteriores.

\begin{figure}[H]
  \centering
  \begin{tikzpicture}[
    font=\footnotesize\sffamily,
    arr/.style={-{Latex}, thick},
    paso/.style={
      draw, thick, rounded corners=2pt, align=center,
      minimum width=3.8cm, minimum height=0.62cm, inner sep=3pt,
      fill=gray!6
    },
    motor/.style={
      draw, thick, rounded corners=2pt, align=center,
      minimum width=3.8cm, minimum height=0.62cm, inner sep=3pt,
      fill=gray!14
    },
    rama/.style={
      draw, thick, rounded corners=2pt, align=center,
      minimum width=3.4cm, minimum height=0.55cm, inner sep=3pt,
      fill=gray!6
    },
    decision/.style={
      diamond, draw, thick, align=center, aspect=2.7,
      inner sep=0.5pt, fill=gray!6, font=\scriptsize\sffamily
    },
    seccion/.style={
      draw, thick, rounded corners=3pt
    },
    seclabel/.style={font=\scriptsize\bfseries\sffamily}
  ]
    % --- Usuario (silueta, sin círculo) ---
    \node[inner sep=0pt, minimum height=0.85cm, minimum width=0.65cm]
      (user) {};
    \begin{scope}[shift={(user.center)}]
      \fill (0,0.26) circle (0.10);
      \draw[thick, line cap=round]
        (0,0.16) -- (0,-0.08)
        (-0.16,0.05) -- (0.16,0.05)
        (0,-0.08) -- (-0.13,-0.32)
        (0,-0.08) -- (0.13,-0.32);
    \end{scope}
    \node[below=0.04cm of user, align=center,
      font=\scriptsize\bfseries\sffamily]
      (userlbl) {Usuario\\[-1pt]{\normalfont\tiny Actor primario}};

    % --- Pestaña Medicamentos ---
    \node[paso, below=0.65cm of userlbl] (selmeds)
      {Seleccionar medicamentos};
    \node[seclabel, above=0.08cm of selmeds] (medslbl)
      {Pestaña: Medicamentos};
    \begin{scope}[on background layer]
      \node[seccion, fit=(medslbl)(selmeds), inner ysep=6pt, inner xsep=9pt]
        (medsbox) {};
    \end{scope}

    % --- Motor y decisión (tras medicamentos) ---
    \node[motor, below=0.5cm of medsbox] (eval)
      {Analizar con el motor};
    \node[decision, below=0.42cm of eval] (salta)
      {¿Salta el\\criterio?};
    \node[rama, below left=0.5cm and 0.02cm of salta] (si)
      {Sí (criterio activo)};
    \node[rama, below right=0.5cm and 0.02cm of salta] (no)
      {No (criterio inactivo)};
    \coordinate (join) at ($(si.south)!0.5!(no.south)$);
    \coordinate (joindown) at ([yshift=-0.48cm]join);

    % --- Pestaña Diagnóstico ---
    \node[paso, below=1.45cm of joindown] (seldx)
      {Seleccionar diagnósticos};
    \node[seclabel, above=0.08cm of seldx] (dxlbl)
      {Pestaña: Diagnóstico};
    \begin{scope}[on background layer]
      \node[seccion, fit=(dxlbl)(seldx), inner ysep=6pt, inner xsep=9pt]
        (dxbox) {};
    \end{scope}

    % --- Motor y decisión (tras diagnósticos) ---
    \node[motor, below=0.5cm of dxbox] (eval2)
      {Analizar con el motor};
    \node[decision, below=0.42cm of eval2] (salta2)
      {¿Salta el\\criterio?};
    \node[rama, below left=0.5cm and 0.02cm of salta2] (si2)
      {Sí (criterio activo)};
    \node[rama, below right=0.5cm and 0.02cm of salta2] (no2)
      {No (criterio inactivo)};
    \coordinate (join2) at ($(si2.south)!0.5!(no2.south)$);
    \coordinate (joindown2) at ([yshift=-0.48cm]join2);

    % --- Resultados ---
    \node[paso, below=1.45cm of joindown2] (res)
      {Resultados finales\\STOPP / START};

    % --- Flechas ---
    \draw[arr] (userlbl.south) -- (medsbox.north);
    \draw[arr] (medsbox.south) -- (eval.north);
    \draw[arr] (eval.south) -- (salta.north);
    \draw[arr] (salta.west) -| node[above, font=\tiny, pos=0.2] {Sí} (si.north);
    \draw[arr] (salta.east) -| node[above, font=\tiny, pos=0.2] {No} (no.north);
    \draw[thick] (si.south) |- (joindown);
    \draw[thick] (no.south) |- (joindown);
    \draw[arr] (joindown) -- (dxbox.north);
    \draw[arr] (dxbox.south) -- (eval2.north);
    \draw[arr] (eval2.south) -- (salta2.north);
    \draw[arr] (salta2.west) -| node[above, font=\tiny, pos=0.2] {Sí} (si2.north);
    \draw[arr] (salta2.east) -| node[above, font=\tiny, pos=0.2] {No} (no2.north);
    \draw[thick] (si2.south) |- (joindown2);
    \draw[thick] (no2.south) |- (joindown2);
    \draw[arr] (joindown2) -- (res.north);
  \end{tikzpicture}
  \caption{Funcionamiento general de la solución: en cada pestaña el motor
  evalúa el caso y marca el criterio como activo o inactivo; el resultado
  final distingue avisos STOPP y START.}
  \label{fig:flujo-vision}
\end{figure}

\section{Metodología de desarrollo}

El desarrollo ha sido iterativo e incremental: cada cambio entra en
incrementos pequeños y verificables, y se contrasta en rondas de revisión
antes de darlo por cerrado. Esta forma de trabajar se justifica por el
dominio del TFG. Un criterio mal codificado puede silenciar un aviso STOPP
o START; conviene detectar ese fallo en el mismo incremento en que se
introduce, contrastando el comportamiento con la guía
oficial~\cite{omahony2023stopp,sacyl2023stopp}.

En la práctica, el trabajo se organizó en rondas temáticas: revisión
sistemática de los criterios STOPP y START por sistemas, una ronda de
corrección del motor y los datos clínicos, y otra de servicios e
interfaz. El detalle de planificación, horas y desviaciones se recoge
en el~\cref{cap:planificacion}; la batería de pruebas, en
el~\cref{cap:pruebas}.

\subsection{Desarrollo dirigido por pruebas (TDD)}

La implementación sigue Test-Driven Development~\cite{beck2002tdd}:
antes de escribir código de producción se formula una prueba que falla
(rojo); a continuación se escribe la implementación mínima que la hace
pasar (verde); después se refactoriza con la suite en verde. El ciclo se
repite en incrementos pequeños, de modo que el sistema permanece en un
estado ejecutable y comprobable.

\begin{figure}[H]
  \centering
  \begin{tikzpicture}[
    font=\small\sffamily,
    caja/.style={
      draw, thick, rounded corners=2pt, align=center,
      minimum width=2.7cm, minimum height=0.9cm, inner sep=3pt,
      fill=gray!6
    },
    arr/.style={-{Latex}, thick}
  ]
    \node[caja] (rojo) {Rojo:\\prueba que falla};
    \node[caja, right=0.85cm of rojo] (verde) {Verde:\\implementación mínima};
    \node[caja, right=0.85cm of verde] (ref) {Refactorizar};
    \draw[arr] (rojo) -- (verde);
    \draw[arr] (verde) -- (ref);
    \draw[arr] (ref.south) -- ++(0,-0.5) -| node[below, font=\scriptsize, pos=0.25]
      {siguiente incremento} (rojo.south);
  \end{tikzpicture}
  \caption{Ciclo rojo--verde--refactor aplicado en el desarrollo.}
  \label{fig:ciclo-tdd}
\end{figure}

En este proyecto TDD no es solo una costumbre de código: es la defensa
principal frente al falso negativo. Si un criterio no salta cuando la
guía dice que debe saltar, el profesional no ve el aviso y el paciente
puede permanecer con una prescripción inapropiada o sin un tratamiento
indicado. Una prueba que describe el caso clínico (medicación,
diagnósticos y, cuando aplica, constantes) y afirma si el criterio
debe activarse o no convierte esa expectativa en una especificación
ejecutable. El mismo esquema cubre el caso contrario ---el criterio no
debe disparar--- para limitar los falsos positivos, que erosionarían la
confianza en la herramienta.

Las pruebas se ejecutan con Karma y Jasmine sobre el código TypeScript
de la aplicación. El criterio de avance de un incremento es que la
nueva prueba pase y que el resto de la suite no se rompa. El inventario
de pruebas y sus resultados se detallan en el~\cref{cap:pruebas}.

\subsection{Control de versiones y flujo de trabajo}

El código se gestiona con Git. Los cambios de entidad ---corrección de
un conjunto de criterios, de un módulo del motor o de la interfaz--- se
desarrollan en ramas de trabajo y se fusionan a la línea principal
cuando la suite de pruebas queda en verde.

Las ideas que alteran el alcance o la interfaz no se implementan
directamente en la línea principal. Permanecen en cuarentena hasta que
hay una decisión explícita de incorporarlas, posponerlas o
descartarlas. Ejemplos de este filtro son la unificación de la
orientación visual entre pestañas, dejada fuera del alcance fusionado, y
el selector jerárquico de variantes diagnósticas, del que solo se
incorporó la primera iteración (familia de hipertensión). De este modo
se separa lo que ya está validado de lo que aún es hipótesis de diseño.
